\documentclass[11pt]{article}
\usepackage[T2A]{fontenc}
\usepackage[utf8]{inputenc}
\usepackage[russian]{babel}
\usepackage[dvipsnames]{xcolor}
\usepackage{tikz}
\usepackage{geometry}
\usepackage{amsmath}
\usepackage{amssymb}
\usepackage{lettrine}
\usepackage{multicol}
\pagestyle{empty}

\geometry{paperwidth=180mm, paperheight=240mm, left=25mm, right=20mm, top=20mm}

\begin{document}
\begin{center}
    48 {\Large КНИГА I ПРЕДЛ. XXIV. ТЕОРЕМА}
\end{center}
\begin{multicols}{2}
\begin{tikzpicture}
    \coordinate (D) at (0, 0);
    \coordinate (A) at (3, 3);
    \coordinate (B) at (4, 0);
    \coordinate (C) at (2, -1);

    \draw[red, dashed] (D) -- (A);
    \draw[red] (A) -- (C);
    \draw[black] (D) -- (B);
    \draw[blue] (D) -- (C);
    \draw[black, dashed] (C) -- (B);
    \draw[blue, dashed] (A) -- (B);

    \node[above] at (A) {A};
    \node[left] at (D) {D};
    \node[below] at (C) {C};
    \node[right] at (B) {B};

    \fill[blue] (D) -- (0:0.5) arc (0:45:0.5);
    \fill[red] (D) -- (-30:0.5) arc (-30:0:0.5);
    \fill[black] (C) -- +(30:0.5) arc (30:80:0.5);
    \fill[yellow] (C) -- +(80:0.5) arc (80:150:0.5);
\end{tikzpicture}

\begin{tikzpicture}
    \coordinate (E) at (3, 4);
    \coordinate (G) at (0, -2);
    \coordinate (F) at (4, 0);

    \node[above] at (E) {E};
    \node[left] at (F) {F};
    \node[below] at (G) {G};

    \draw[yellow] (G) -- (F);
    \draw[red] (G) -- (E);
    \draw[blue] (E) -- (F);

    \draw[black, dashed] (E) -- +(240:0.8) arc (240:290:0.8);
\end{tikzpicture}
\columnbreak \\
Если у двух треугольников по две стороны со-
ответственно равны друг другу (
\begin{tikzpicture}
    \draw[blue] (0,0) node[left]{A} -- (1, 0) node[right]{B};
\end{tikzpicture}
=
\begin{tikzpicture}
    \draw[blue] (0,0) node[left]{E} -- (1, 0) node[right]{F};
\end{tikzpicture}
и
\begin{tikzpicture}
    \draw[red] (0,0) node[left]{A} -- (1, 0) node[right]{D};
\end{tikzpicture}
=
\begin{tikzpicture}
    \draw[red, dashed] (0,0) node[left]{G} -- (1, 0) node[right]{E};
\end{tikzpicture}
), и угол заключенный ними в одном
\begin{tikzpicture}
    \draw[red, dashed] (0,0) node[left]{B} -- (0.7, 1) node[right]{A};
    \draw[red] (0.7, 1) -- (0.7, 0);
    \draw[blue] (0.7, 1) -- (1.4, 0) node[right]{B};
\end{tikzpicture}
больше, чем в другом
\begin{tikzpicture}
    \draw[red] (0,0) node[left]{E} -- (0.3, 1) node[right]{F};
    \draw[blue] (0.3, 1) -- (0.6, 0) node[right]{G};
\end{tikzpicture}
, то сторона
\begin{tikzpicture}
    \draw[black] (0, 0) node[left]{D} -- (1, 0) node[right]{B};
\end{tikzpicture}
противолежащая большему углу больше стороны, противолежащей меньшему
\begin{tikzpicture}
    \draw[yellow] (0, 0) node[left]{F} -- (1, 0) node[right]{G}
\end{tikzpicture}
\\
Сделаем
\begin{tikzpicture}
    \draw[red] (0,0) node[left]{C} -- (0.5, 1) node[right]{A};
    \draw[blue] (0.5, 1) -- (1, 0) node[right]{B};
\end{tikzpicture}
=
\begin{tikzpicture}
    \draw[red] (0,0) node[left]{G} -- (0.5, 1) node[right]{E};
    \draw[blue] (0.5, 1) -- (1, 0) node[right]{F};
\end{tikzpicture}
и
\begin{tikzpicture}
    \draw[red] (0,0) node[left]{C} -- (1, 0) node[right]{A};
\end{tikzpicture}
=
\begin{tikzpicture}
    \draw[red] (0,0) node[left]{G} -- (1, 0) node[right]{E};
\end{tikzpicture}
\\
проведём
\begin{tikzpicture}
    \draw[blue, dashed] (0,0) node[left]{C} -- (1, 0) node[right]{D};
\end{tikzpicture}
и
\begin{tikzpicture}
    \draw[black, dashed] (0,0) node[left]{B} -- (1, 0) node[right]{C};
\end{tikzpicture}
Поскольку
\begin{tikzpicture}
    \draw[red] (0,0) node[left]{C} -- (1, 0) node[right]{A};
\end{tikzpicture}
=
\begin{tikzpicture}
    \draw[red, dashed] (0,0) node[left]{A} -- (1, 0) node[right]{D};
\end{tikzpicture}
\begin{tikzpicture}
    \fill[black] (0,0) circle (1pt);
    \fill[black] (0.1,0.1) circle (1pt);
    \fill[black] (0.2,0) circle (1pt);
\end{tikzpicture}

\begin{tikzpicture}
    \coordinate (D) at (0, 0);
    \coordinate (A) at (0.5, 0.5);
    \coordinate (C) at (0.5, -0.5);

    \node[above] at (A) {A};
    \node[left] at (D) {D};
    \node[below] at (C) {C};
    
    \fill[blue] (D) -- +(0:0.5) arc (0:45:0.5);
    \fill[red] (D) -- +(-45:0.5) arc (-45:0:0.5);
\end{tikzpicture}
{\Large =}
\begin{tikzpicture}
    \coordinate (D) at (0, 0);
    \coordinate (A) at (0.5, 0.5);
    \coordinate (C) at (0.5, -0.5);

    \node[above] at (A) {A};
    \node[left] at (D) {D};
    \node[below] at (C) {C};
    
    \fill[yellow] (D) -- +(-45:0.5) arc (-45:45:0.5);
\end{tikzpicture}
,но\\
\begin{tikzpicture}
    \coordinate (D) at (0, 0);
    \coordinate (B) at (0.5, 0);
    \coordinate (C) at (0.5, -0.5);

    \node[above] at (B) {B};
    \node[left] at (D) {D};
    \node[below] at (C) {C};
    
    \fill[red] (D) -- +(-45:0.5) arc (-45:0:0.5);
\end{tikzpicture}
{\Large <}
\begin{tikzpicture}
    \coordinate (D) at (0, 0);
    \coordinate (A) at (0.5, 0.5);
    \coordinate (C) at (0.5, -0.5);

    \node[above] at (A) {A};
    \node[left] at (D) {D};
    \node[below] at (C) {C};
    
    \fill[yellow] (D) -- +(-45:0.5) arc (-45:45:0.5);
\end{tikzpicture}
И
\begin{tikzpicture}
    \fill[black] (0,0) circle (1pt);
    \fill[black] (0.1,0.1) circle (1pt);
    \fill[black] (0.2,0) circle (1pt);
\end{tikzpicture}

\begin{tikzpicture}
    \draw[black] (0,0) node[left]{D} -- (1, 0) node[right]{B};
\end{tikzpicture}
=
\begin{tikzpicture}
    \draw[black, dashed] (0,0) node[left]{B} -- (1, 0) node[right]{C};
\end{tikzpicture}
,но
\begin{tikzpicture}
    \draw[black, dashed] (0,0) node[left]{B} -- (1, 0) node[right]{C};
\end{tikzpicture}
=
\begin{tikzpicture}
    \draw[yellow] (0,0) node[left]{F} -- (1, 0) node[right]{G};
\end{tikzpicture}

\begin{tikzpicture}
    \fill[black] (0,0) circle (1pt);
    \fill[black] (0.1,0.1) circle (1pt);
    \fill[black] (0.2,0) circle (1pt);
\end{tikzpicture}
\begin{tikzpicture}
    \draw[black] (0,0) node[left]{D} -- (1, 0) node[right]{B};
\end{tikzpicture}
>
\begin{tikzpicture}
    \draw[yellow] (0,0) node[left]{F} -- (1, 0) node[right]{G};
\end{tikzpicture}
\\
ч. т. д.
\end{multicols}
\end{document}